% [Poster title]: conference poster, A0 portrait.
% Laid out on a 1000 x 1414.2 mm design page (A0 proportions) and placed on the 841 x 1189 mm sheet by
% poster_A0.tex, so every size here prints at 0.841 of its design size (88 pt title -> 74 pt, 27 pt body -> 22.7 pt).
%
% PROVENANCE (keep this block current; it is the record of what came from where):
%   figures/[fig_overview.pdf]  = the paper's Fig. [1], unchanged   [or: cropped to panels a-c; text outlined]
%   figures/[fig_method.pdf]    = the paper's Fig. [2]
%   tables/results.tex          = the manuscript's table source, unchanged except [comparator citations added]
%   Prose and maths: pasted from the paper ([main text, appendix]), shortened only by cutting.
%   Departures from the paper's wording: [acronyms spelled out at first use; plain-language table key].
%
% Layout: header; four bands (Abstract | Key idea; Method; Main result and gallery; Data | Results); footer.
% Every piece is a TikZ node chained to the one above it, so nothing can overlap; the \foreach at the end logs
% every box for box_check.py. Missing figure files show as grey boxes labelled with their path.
% Build: make (two xelatex passes, because absolute placement on the page needs the second pass).
\documentclass{article}
\usepackage[paperwidth=1000mm, paperheight=1414.2mm, margin=0mm]{geometry}   % A0 aspect; printed at 0.841
\usepackage{fontspec}
\IfFontExistsTF{Helvetica Neue}{%
  \setmainfont{Helvetica Neue}[UprightFont=* Light, BoldFont=* Medium, ItalicFont=* Light Italic,
    BoldItalicFont=* Medium Italic]
  \newfontfamily\heavy{Helvetica Neue Bold}
  \newfontfamily\thin{Helvetica Neue Thin}
}{% fallback where Helvetica Neue is not installed (Suggested, not yet tested in real work)
  \IfFontExistsTF{texgyreheros-regular.otf}{\setmainfont{texgyreheros}[Extension=.otf, UprightFont=*-regular,
    BoldFont=*-bold, ItalicFont=*-italic, BoldItalicFont=*-bolditalic]}{\setmainfont{lmsans10}[Extension=.otf,
    UprightFont=*-regular, BoldFont=*-bold, ItalicFont=*-oblique, BoldItalicFont=*-boldoblique]}
  \newcommand{\heavy}{\bfseries}\newcommand{\thin}{}
}
\usepackage{amsmath}
\usepackage[mathrm=sym]{unicode-math}
% maths is set as in the paper: Computer Modern (Latin Modern Math, also for \mathrm), at the text's x-height
\setmathfont{latinmodern-math.otf}[Scale=MatchLowercase]
\usepackage{microtype}
\usepackage{graphicx}
\usepackage[table]{xcolor}
\usepackage{booktabs,array,tabularx}
\usepackage{tikz}
\usetikzlibrary{calc,positioning,fit}
\pagestyle{empty}
\setlength{\parindent}{0pt}

% ---------- poster chrome (figures keep the paper's own colors) ----------
\definecolor{ink}{HTML}{1A1A1A}
\definecolor{slate}{HTML}{4A4F5A}
\definecolor{accent}{HTML}{2B50AA}   % the one accent: the method's name, the coined term, key phrases. Use the
                                     % host's or the paper's color here.

% ---------- type scale (design pt; printed = x 0.841) ----------
\newcommand{\sectfont}{\fontsize{46}{54}\selectfont\bfseries}
\newcommand{\absfont}{\fontsize{33}{42.5}\selectfont}
\newcommand{\bodyfont}{\fontsize{27}{35}\selectfont}
\newcommand{\captfont}{\fontsize{25}{31}\selectfont}
\newcommand{\mathfont}{\fontsize{27}{34}\selectfont}

% ---------- a figure of the paper, or a labelled grey box while the file is missing ----------
% \fig[includegraphics options]{width cm}{placeholder height cm}{file}
\newcommand{\fig}[4][]{\IfFileExists{#4}{\includegraphics[width=#2cm,#1]{#4}}{%
  {\setlength{\fboxsep}{0pt}\setlength{\fboxrule}{0.06cm}\fcolorbox{slate}{black!4}{%
  \begin{minipage}[c][#3cm][c]{\dimexpr#2cm-0.12cm\relax}\centering\captfont\color{slate}%
  [figure: \detokenize{#4}]\end{minipage}}}}}

% ---------- the manuscript's table, unchanged, as a plain box ----------
% \resultstable{width}: for a table that sets itself to \linewidth (tabularx, or \resizebox{\linewidth}{!}{...})
% \resultstablescaled{width}: for a plain tabular at its natural width, scaled as a whole to the strip's width
\newcommand{\resultstable}[1]{\begingroup
  \renewenvironment{table}[1][]{}{}\renewcommand{\caption}[2][]{}\renewcommand{\label}[1]{}%
  \begin{minipage}{#1}% PLACEHOLDER: replace this whole file with the manuscript's own table source, copied unchanged.
% The poster's \resultstable wrapper neutralizes table, \caption and \label, so the manuscript file goes in as it is.
% The only edit allowed: the comparator's citation number after its name, matching the poster's reference list.
% Every value below is a placeholder in [brackets]; none is a result. A cell that starts with [ right after \\
% or a rule must be braced, {[...]}, or LaTeX reads it as an optional argument.
\begin{table}[t]
\centering
\definecolor{rankgreen}{HTML}{159A6C}
\definecolor{rankblue}{HTML}{0A84FF}
\definecolor{rankorange}{HTML}{F5A20B}
\definecolor{soft}{RGB}{105,105,105}
\caption{[The manuscript's caption: it is dropped on the poster.]}
\label{tab:results}
\begingroup
\fontsize{24}{29}\selectfont
\setlength{\tabcolsep}{6pt}
\renewcommand{\arraystretch}{1.15}
\newcommand{\rfirst}[2]{\cellcolor{rankgreen!19}\textbf{#1}{\fontsize{17}{19}\selectfont\color{soft}$\pm$#2}}
\newcommand{\rsecond}[2]{\cellcolor{rankblue!16}\textbf{#1}{\fontsize{17}{19}\selectfont\color{soft}$\pm$#2}}
\newcommand{\rthird}[2]{\cellcolor{rankorange!20}\textbf{#1}{\fontsize{17}{19}\selectfont\color{soft}$\pm$#2}}
\newcommand{\res}[2]{#1{\fontsize{17}{19}\selectfont\color{soft}$\pm$#2}}
\begin{tabularx}{\linewidth}{@{}l>{\raggedright\arraybackslash}X cccc c@{}}
\toprule
\textbf{method} & \textbf{setting} & \textbf{[task 1]} & \textbf{[task 2]} & \textbf{[task 3]} & \textbf{[task 4]}
  & \textbf{mean rank} $\downarrow$\\
\midrule
\multicolumn{7}{@{}l}{\textit{[group of comparators]}}\\
{[Comparator A]}\,\cite{5} & [setting] & \rsecond{[.00]}{[.00]} & \res{[.00]}{[.00]} & \rthird{[.00]}{[.00]} & \res{[.00]}{[.00]} & [0.0]\\
{[Comparator B]} & [setting] & \res{[.00]}{[.00]} & \rthird{[.00]}{[.00]} & \res{[.00]}{[.00]} & \rsecond{[.00]}{[.00]} & [0.0]\\
{[Comparator C]} & [setting] & \rthird{[.00]}{[.00]} & \res{[.00]}{[.00]} & \rsecond{[.00]}{[.00]} & \res{[.00]}{[.00]} & [0.0]\\
\multicolumn{7}{@{}l}{\textit{[baselines]}}\\
{[Baseline]} & [setting] & \res{[.00]}{[.00]} & \rsecond{[.00]}{[.00]} & \res{[.00]}{[.00]} & \rthird{[.00]}{[.00]} & [0.0]\\
\midrule
\textbf{[Method] (ours)} & [setting] & \rfirst{[.00]}{[.00]} & \rfirst{[.00]}{[.00]} & \rfirst{[.00]}{[.00]} & \rfirst{[.00]}{[.00]} & \textbf{[0.0]}\\
\bottomrule
\end{tabularx}
\endgroup
\end{table}
\end{minipage}\endgroup}
\newsavebox{\tablebox}
\newcommand{\resultstablescaled}[1]{\begingroup
  \renewenvironment{table}[1][]{}{}\renewcommand{\caption}[2][]{}\renewcommand{\label}[1]{}%
  \renewcommand{\centering}{}\sbox{\tablebox}{% PLACEHOLDER: replace this whole file with the manuscript's own table source, copied unchanged.
% The poster's \resultstable wrapper neutralizes table, \caption and \label, so the manuscript file goes in as it is.
% The only edit allowed: the comparator's citation number after its name, matching the poster's reference list.
% Every value below is a placeholder in [brackets]; none is a result. A cell that starts with [ right after \\
% or a rule must be braced, {[...]}, or LaTeX reads it as an optional argument.
\begin{table}[t]
\centering
\definecolor{rankgreen}{HTML}{159A6C}
\definecolor{rankblue}{HTML}{0A84FF}
\definecolor{rankorange}{HTML}{F5A20B}
\definecolor{soft}{RGB}{105,105,105}
\caption{[The manuscript's caption: it is dropped on the poster.]}
\label{tab:results}
\begingroup
\fontsize{24}{29}\selectfont
\setlength{\tabcolsep}{6pt}
\renewcommand{\arraystretch}{1.15}
\newcommand{\rfirst}[2]{\cellcolor{rankgreen!19}\textbf{#1}{\fontsize{17}{19}\selectfont\color{soft}$\pm$#2}}
\newcommand{\rsecond}[2]{\cellcolor{rankblue!16}\textbf{#1}{\fontsize{17}{19}\selectfont\color{soft}$\pm$#2}}
\newcommand{\rthird}[2]{\cellcolor{rankorange!20}\textbf{#1}{\fontsize{17}{19}\selectfont\color{soft}$\pm$#2}}
\newcommand{\res}[2]{#1{\fontsize{17}{19}\selectfont\color{soft}$\pm$#2}}
\begin{tabularx}{\linewidth}{@{}l>{\raggedright\arraybackslash}X cccc c@{}}
\toprule
\textbf{method} & \textbf{setting} & \textbf{[task 1]} & \textbf{[task 2]} & \textbf{[task 3]} & \textbf{[task 4]}
  & \textbf{mean rank} $\downarrow$\\
\midrule
\multicolumn{7}{@{}l}{\textit{[group of comparators]}}\\
{[Comparator A]}\,\cite{5} & [setting] & \rsecond{[.00]}{[.00]} & \res{[.00]}{[.00]} & \rthird{[.00]}{[.00]} & \res{[.00]}{[.00]} & [0.0]\\
{[Comparator B]} & [setting] & \res{[.00]}{[.00]} & \rthird{[.00]}{[.00]} & \res{[.00]}{[.00]} & \rsecond{[.00]}{[.00]} & [0.0]\\
{[Comparator C]} & [setting] & \rthird{[.00]}{[.00]} & \res{[.00]}{[.00]} & \rsecond{[.00]}{[.00]} & \res{[.00]}{[.00]} & [0.0]\\
\multicolumn{7}{@{}l}{\textit{[baselines]}}\\
{[Baseline]} & [setting] & \res{[.00]}{[.00]} & \rsecond{[.00]}{[.00]} & \res{[.00]}{[.00]} & \rthird{[.00]}{[.00]} & [0.0]\\
\midrule
\textbf{[Method] (ours)} & [setting] & \rfirst{[.00]}{[.00]} & \rfirst{[.00]}{[.00]} & \rfirst{[.00]}{[.00]} & \rfirst{[.00]}{[.00]} & \textbf{[0.0]}\\
\bottomrule
\end{tabularx}
\endgroup
\end{table}
}\resizebox{#1}{!}{\usebox{\tablebox}}\endgroup}

% ---------- grammar: section title + rule, then pieces chained below ----------
% Text is fully justified. TikZ sets \spaceskip from the 10 pt font before the font switch inside text-width
% nodes; \txt resets it so word spaces come from the poster font. \hyphenpenalty 1000: 5000 left loose lines
% and one-word last lines.
\tikzset{
  piece/.style={anchor=north west, inner sep=0pt, outer sep=0pt},
  sect/.style={piece, font=\sectfont, text=ink, text height=1.22cm, text depth=0.30cm},
}
\newcommand{\rulebelow}[2]{\fill[ink] ($(#1.south west)+(0,-0.3)$) rectangle ++(#2,-0.09);}
\newcommand{\txt}[3]{\begin{minipage}[t]{\linewidth}#1\spaceskip=0pt\xspaceskip=0pt\frenchspacing
  \hyphenpenalty=1000\exhyphenpenalty=1000\emergencystretch=2em\color{#2}#3\end{minipage}}
\newcommand{\body}[1]{\txt{\bodyfont}{ink}{#1}}
\newcommand{\capt}[1]{\txt{\captfont}{slate}{#1}}
\newcommand{\hi}[1]{\textbf{\color{accent}#1}}
\renewcommand{\cite}[1]{{\color{slate}[#1]}}
\newcommand{\bignum}[1]{{\thin\fontsize{72}{72}\selectfont\color{accent}#1}}
\newcommand{\numlab}[1]{{\captfont\color{slate}#1}}
% label under a brace: every label gets one height and depth, so all labels share a baseline
\newcommand{\lab}[1]{\text{\fontsize{22}{26}\selectfont\color{slate}\vphantom{fp}#1}}
\newcommand{\datatable}[1]{\begingroup\fontsize{27}{32}\selectfont\spaceskip=0pt\renewcommand{\arraystretch}{1.28}%
  \setlength{\heavyrulewidth}{2pt}\setlength{\lightrulewidth}{1.1pt}\setlength{\aboverulesep}{0.18em}%
  \setlength{\belowrulesep}{0.22em}#1\endgroup}

\begin{document}
\begin{tikzpicture}[remember picture, overlay, x=1cm, y=1cm]
\coordinate (O) at (current page.south west);
\newcommand{\at}[2]{($(O)+(#1,#2)$)}

% =====================================================================
% HEADER (black on white): title, authors, affiliations, a full-width rule
% =====================================================================
\node[piece] (title) at \at{3}{139.32} {\begin{minipage}[t]{72cm}\raggedright
  \fontsize{88}{98}\selectfont\heavy\textcolor{accent}{[Method]}: [the paper's title,\\as printed on the paper]\end{minipage}};
% host logo only if the venue or the host asks for one; delete this node (and its name in the \foreach) otherwise
\node[anchor=north east, inner sep=0pt] (logo) at \at{97}{139.27} {\fig{21}{6}{figures/[logo.pdf]}};
\node[piece, below=0.6cm of title.south west, anchor=north west] (authors) {\fontsize{35}{42}\selectfont
  \underline{[Presenting Author]}\textsuperscript{1}, [Second Author]\textsuperscript{2},
  [Third Author]\textsuperscript{1,\dag}, [Last Author]\textsuperscript{2,\dag}};
\node[piece, below=0.45cm of authors.south west, anchor=north west, text width=94cm] (affil) {%
  \begin{minipage}[t]{94cm}\raggedright\fontsize{23}{29}\selectfont\spaceskip=0pt\xspaceskip=0pt\color{slate}%
  \mbox{\textsuperscript{1}[Department, Institution, City]}\quad
  \mbox{\textsuperscript{2}[Department, Institution, City]}\quad
  \mbox{\textsuperscript{\dag}These authors jointly supervised this work.}\end{minipage}};
\coordinate (hdr) at ($(affil.south west)+(0,-0.55)$);
\fill[ink] (hdr) rectangle ($(hdr)+(94,-0.2)$);

% =====================================================================
% BAND 1: Abstract (55.5 cm) | Key idea: the paper's concept figure and its defining equations (36.5 cm)
% =====================================================================
\def\Aw{55.5}
\node[sect, below=1.0cm of hdr, anchor=north west] (s1) {Abstract};
\rulebelow{s1}{\Aw}
\node[piece, below=0.75cm of s1.south west, anchor=north west, text width=\Aw cm] (abs) {\txt{\absfont}{ink}{%
  \textbf{[The abstract's opening question or claim, pasted.]} [Paste the paper's abstract here, sentence by
  sentence, and shorten it only by cutting whole sentences or clauses. Do not paraphrase, merge two sentences or
  add a connective. Mark four or five central phrases with] \hi{[the accent macro]}[, and the method's name in]
  \textbf{[bold]}. [This placeholder runs to about the length that fills the column at this size: the abstract
  of the poster it was taken from held 170 words at 33 pt design size. If the paper's abstract is longer, cut
  the sentences that restate the setting, keep the question, the idea, what it makes possible and the one
  number that carries the result. Spell out an acronym at its first use on the poster and record that departure
  in the provenance block at the top of this file, because the poster has no glossary and a reader at a poster
  session does not read the introduction first.]}};

\def\Bx{60.5}\def\Bw{36.5}
\node[sect] (s2) at ($(s1.north west)+(\Bx-3,0)$) {[Key idea]};
\rulebelow{s2}{\Bw}
\node[piece, below=0.75cm of s2.south west, anchor=north west] (sph)
  {\fig[trim=0 0 0 0, clip]{15.6}{15.6}{figures/[fig_concept.pdf]}};
\node[piece, text width=19.4cm] (defs) at ($(sph.north east)+(1.5,0)$) {%
  \begin{minipage}[t]{19.4cm}\raggedright\mathfont\spaceskip=0pt\color{ink}%
  {\captfont\color{slate}[label of the first definition]}\\[0.05em]
  $[\,\text{first defining equation of the paper}\,]$\\[0.3em]
  {\setlength{\fboxrule}{2.2pt}\setlength{\fboxsep}{0.27cm}\fcolorbox{accent}{accent!7}{%
    \begin{minipage}{\dimexpr19.4cm-2\fboxrule-2\fboxsep\relax}\centering
    {\fontsize{27}{32}\selectfont\bfseries\color{accent}[name of the core definition]}%
    \fontsize{28}{34}\selectfont\abovedisplayskip=0.3em\belowdisplayskip=0.4em
    \abovedisplayshortskip=\abovedisplayskip\belowdisplayshortskip=\belowdisplayskip\jot=0.35em
    \begin{gather*}
      [\,\text{the core definition, first line}\,]\\
      [\,\text{the core definition, second line}\,]
    \end{gather*}
    {\captfont\color{slate}[one sentence of the paper about it]}\end{minipage}}}\\[0.3em]
  {\captfont\color{slate}[label of a second property]}\\[0.05em]
  $[\,\text{second equation}\,]$\end{minipage}};

% =====================================================================
% BAND 2: Method, the paper's two method figures side by side (2 x 45.85 cm, gutter 2.3 cm)
% =====================================================================
\node[fit=(abs)(sph)(defs), inner sep=0pt] (b1) {};
\node[sect] (s3) at ($(b1.south west)+(0,-0.86)$) {[Method]};
\rulebelow{s3}{94}
\def\Fw{45.85}
\node[piece, below=0.75cm of s3.south west, anchor=north west] (f1) {\fig{\Fw}{22}{figures/[fig_overview.pdf]}};
\node[piece] (f2) at ($(f1.north west)+(48.15,0)$) {\fig{\Fw}{22}{figures/[fig_method.pdf]}};
\node[fit=(f1)(f2), inner sep=0pt] (f12) {};
\node[piece, text width=\Fw cm] (c3) at ($(f1.south west |- f12.south)+(0,-0.6)$) {\capt{%
  {\centering [The paper's sentence that names the training objective]\,\cite{1,2}:\par}\vspace{0.15em}
  {\centering\mathfont\color{ink}$\displaystyle\mathcal L=\underbrace{\mathcal L_{\mathrm{a}}}_{\lab{[first term]}}
  +\underbrace{\lambda\mathcal L_{\mathrm{b}}\vphantom{\mathcal L_{\mathrm{a}}}}_{\lab{[second term]}}$\par}}};
\node[piece, text width=\Fw cm] (c4) at ($(f2.south west |- f12.south)+(0,-0.6)$) {\capt{%
  [The paper's caption sentences for this figure, cut. Name each component as the paper does and cite its
  source]\,\cite{3}. [An acronym is spelled out at first use, for example Long Short-Term Memory
  (LSTM)]\,\cite{4}.}};

% =====================================================================
% BAND 3: the main result figure with its text column, then a two-figure gallery
% =====================================================================
\node[fit=(c3)(c4), inner sep=0pt] (b2) {};
\node[sect] (s4) at ($(b2.south west)+(0,-0.86)$) {[Main result section title]};
\rulebelow{s4}{94}
\def\Wa{66}\pgfmathsetmacro\Ta{92-\Wa}
\node[piece, below=0.75cm of s4.south west, anchor=north west] (fa) {\fig{\Wa}{12.5}{figures/[fig_result.pdf]}};
\node[piece, text width=\Ta cm] (ca) at ($(fa.north east)+(2,0)$) {\txt{\bodyfont}{slate}{%
  \textbf{\color{ink}[The figure's bold caption claim.]} [The caption's next sentences, cut: what is shown,
  at what scale and what the colors of the figure encode.]}};
\node[piece, below=0.6cm of ca.south west, anchor=north west, text width=\Ta cm] (ca2) {\body{%
  [One or two sentences of the results text, with] \textbf{[the key phrase]} [in bold.]}};
\node[fit=(fa)(ca)(ca2), inner sep=0pt] (fac) {};
\node[piece, below=1.0cm of fac.south west, anchor=north west, text width=45.85cm] (l5a) {\capt{%
  \textbf{\color{ink}[Bold caption claim of the left gallery figure.]}}};
\node[piece, text width=45.85cm] (l5b) at ($(l5a.north west)+(48.15,0)$) {\capt{%
  \textbf{\color{ink}[Bold caption claim of the right gallery figure.]}}};
\node[piece, below=0.45cm of l5a.south west, anchor=north west] (g1) {\fig{45.85}{24}{figures/[fig_gallery_a.pdf]}};
\node[piece] (g2) at ($(g1.north west)+(48.15,0)$) {\fig{45.85}{24}{figures/[fig_gallery_b.pdf]}};

% =====================================================================
% BAND 4: Data (31 cm): big numbers and the inventory | Results (61 cm): the manuscript's table and its key
% =====================================================================
\node[fit=(g1)(g2), inner sep=0pt] (b3) {};
\def\Lw{31}\pgfmathsetmacro\Rx{\Lw+5}\pgfmathsetmacro\Rw{92-\Lw}
\node[sect] (s6) at ($(b3.south west)+(0,-0.86)$) {[Data]};
\rulebelow{s6}{\Lw}
\node[piece, below=0.75cm of s6.south west, anchor=north west] (n6) {%
  \begin{tabular*}{\Lw cm}{@{\extracolsep{\fill}}lll@{}}
  \bignum{[0.0\,B]} & \bignum{[0.0\,M]} & \bignum{[000]}\\[0.3cm]
  \numlab{[unit one]} & \numlab{[unit two]} & \numlab{[unit three]}
  \end{tabular*}};
% the paper's own inventory totals; B = billion, M = million (write the word in the label if the room is mixed)
\node[piece, below=0.7cm of n6.south west, anchor=north west] (t6) {\datatable{%
  \begin{tabular*}{\Lw cm}{@{\extracolsep{\fill}}lrr@{}}
  \toprule
  \textbf{[group]} & \textbf{[count]} & \textbf{[size]}\\
  \midrule
  {[group one]} & [000] & [0.0\,M]\\
  {[group two]} & [000] & [0.0\,M]\\
  {[further groups]} & [000] & [0.0\,M]\\
  \midrule
  \textbf{total} & \textbf{[000]} & \textbf{[0.0\,M]}\\
  \bottomrule
  \end{tabular*}}};

\node[sect, below=1.3cm of t6.south west, anchor=north west] (s7) {[Evaluation]};
\rulebelow{s7}{\Lw}
\node[piece, below=0.75cm of s7.south west, anchor=north west, text width=\Lw cm] (c7) {\body{%
  [The paper's sentence on the benchmark: how many held-out sets, how many tasks.] [The protocol in one
  sentence, with] \textbf{[the key phrase]}[.]}};

\node[sect] (s9) at ($(s6.north west)+(\Rx-3,0)$) {[Results]};
\node[piece, below=0.23cm of s9.south west, anchor=north west] (t9) {\resultstable{\Rw cm}};
\node[piece, below=0.4cm of t9.south west, anchor=north west, text width=\Rw cm] (c9) {\body{%
  [The paper's one-sentence takeaway of this table, pasted, with] \textbf{[its key phrase]} [in bold.]}};
\node[piece, below=0.35cm of c9.south west, anchor=north west, text width=\Rw cm] (c9b) {\capt{%
  [The table key: what each metric is and which direction is better; what the shading and bold mean; what the
  error after each value is.]}};

% =====================================================================
% FOOTER (black on white): a full-width rule, references and contacts on the left, the release line on the right
% =====================================================================
\newcommand{\rsep}{\hspace{0.4em plus 0.2em minus 0.05em}}   % between references
\def\Fw{77.4}
\node[anchor=south east, inner sep=0pt] (npt) at \at{97}{1.3} {\fontsize{26}{31}\selectfont\bfseries
  [Preprint and code coming soon]};
\node[anchor=south, inner sep=0pt] (np) at ($(npt.north)+(0,0.35)$) {\fig{2.4}{2.4}{assets/no_photos.pdf}};
\node[anchor=south west, inner sep=0pt] (contacts) at \at{3}{1.3} {\begin{minipage}{\Fw cm}%
  \fontsize{22}{26}\selectfont\spaceskip=0pt\color{slate}%
  [you@example.org]\enspace[social handle]\hfill
  [lab name]\enspace[lab web address]\hfill
  [second lab]\enspace[web address]\end{minipage}};
\node[anchor=south west, inner sep=0pt] (refs) at ($(contacts.north west)+(0,0.4)$) {\begin{minipage}{\Fw cm}%
  \fontsize{18}{22.5}\selectfont\spaceskip=0pt\frenchspacing\color{slate}\textbf{\color{ink}References}\rsep
  [1]~[First author]~\emph{et~al.}, [short title or method name], [Venue]~[Year]\rsep
  [2]~[First author]~\emph{et~al.}, [short title], [Venue]~[Year]\rsep
  [3]~[First author]~\emph{et~al.}, [short title], [Venue]~[Year]\rsep
  [4]~[First author]~\emph{et~al.}, [short title], [Venue]~[Year]\rsep
  [5]~[First author]~\emph{et~al.}, [comparator in the table], [Venue]~[Year]\rsep
  Data: [First author]~\emph{et~al.}, [dataset], [Venue]~[Year]\end{minipage}};
\node[fit=(refs)(contacts)(np)(npt), inner sep=0pt] (foot) {};
\fill[ink] ($(O |- foot.north)+(3,0.55)$) rectangle ($(O |- foot.north)+(97,0.75)$);
\node[inner sep=0pt, minimum height=0.2cm, minimum width=94cm, anchor=south west] (footrule) at ($(O |- foot.north)+(3,0.55)$) {};

% ---------- log every box in page cm for box_check.py (add every new node here) ----------
\foreach \nm in {title,logo,authors,affil,s1,abs,s2,sph,defs,s3,f1,f2,c3,c4,s4,fa,ca,ca2,l5a,l5b,g1,g2,s6,n6,t6,s7,c7,
                 s9,t9,c9,c9b,np,npt,refs,contacts,footrule}{
  \path let \p0=(O), \p1=(\nm.south west), \p2=(\nm.north east) in
    \pgfextra{\pgfmathsetmacro\xa{(\x1-\x0)/28.45276}\pgfmathsetmacro\ya{(\y1-\y0)/28.45276}%
      \pgfmathsetmacro\xb{(\x2-\x0)/28.45276}\pgfmathsetmacro\yb{(\y2-\y0)/28.45276}%
      \typeout{BOX \nm\space\xa\space\ya\space\xb\space\yb}};}
\end{tikzpicture}
\end{document}
